\documentclass[../main.tex]{subfiles}

\begin{document}
% -----------------------------
\section{Harness}
% -----------------------------
We can also create a small a small driver program that invokes compiled code and allows
it to be executed and tested inside a normal application.
We implement the harness in C because the compiled function follows the C calling convention,
making it straightforward to import it with extern and link it directly into the executable.

Create a file \path{back_end/harness.c}:
\begin{lstlisting} [language=C, caption={Simple C application}]
extern int my_func(int);

#include <stdio.h>
#include <stdlib.h>

int main(int argc, char **argv) {
    int a = atoi(argv[1]);   // convert command-line string to int
    printf("Result of execution for (%d): %d\n", a, my_func(a));
}
\end{lstlisting}
And extend \path{back_end.sh}:
\begin{lstlisting} [language=Python, caption={}]
clang ${SCRIPT_DIR}/harness.c ${BUILD_DIR}/${MODEL_STEM}_out.o -o \
    ${BUILD_DIR}/${MODEL_STEM}_run

echo "Created executable: ${BUILD_DIR}/${MODEL_STEM}_run"
echo "Running executable..."

${BUILD_DIR}/${MODEL_STEM}_run 42
\end{lstlisting}

By running our function with different values, we can see that results will produce
a U-shaped curve with a minimum at 40, like expected.
\begin{figure}[h]
    \centering
    \includegraphics[width=0.5\textwidth]{run_results.png}
    \caption{Execution results}
    \label{fig:results}
\end{figure}

\end{document}
